\section{Theory proofs, extensions, and scope}
\label{app:theory}

This appendix gives complete proofs for Section~\ref{sec:theory}.  It also
records edge cases that are easy to lose when a soft router is visualized as a
discrete sharing graph.

\subsection{Static SSMF gauge details}
\label{app:gauge_proof}

The non-identifiability conclusion and the symmetric uniformizing path are
the unconstrained SSMF result of
\citet[Section~IV-A]{vuthanh2023bounded} under
\(X=\Theta^\top\), \(W=B^\top\), and \(H=A^\top\); see
Appendix~\ref{app:literature}.  The proof below only records the direct local
row-sum-preserving extension used to delimit which router summaries change.

\begin{proof}[Proof of Proposition~\ref{thm:continuous_gauge}]
The constraint $H\ones=0$ implies
\begin{equation}
  M_t\ones=(I+tH)\ones=\ones.
  \label{eq:app_rowsum}
\end{equation}
The determinant is continuous and $\det(M_0)=1$, so $M_t$ is invertible for
all sufficiently small $|t|$.  Strict positivity is also an open condition:
because $A$ has finitely many entries and $A_{ij}>0$, there is a neighborhood
of zero in which every entry of $AM_t$ remains positive.  Therefore
\begin{equation}
  A_t\ones=AM_t\ones=A\ones=\ones,
  \qquad A_t>0,
\end{equation}
and $A_t\in\simplex_K^L$.  The effective parameters are unchanged because
\begin{equation}
  A_tB_t=(AM_t)(M_t^{-1}B)=AB.
\end{equation}

\Eqref{eq:app_rowsum} also gives
$M_t^{-1}\ones=\ones$.  Hence the augmented basis matrix satisfies
\begin{equation}
 [B_t,\ones]
 =[M_t^{-1}B,M_t^{-1}\ones]
 =M_t^{-1}[B,\ones].
 \label{eq:app_affine_rank}
\end{equation}
Left multiplication by an invertible matrix preserves row rank, so affine
independence is preserved exactly, not merely infinitesimally.

It remains to rule out a permutation.  If $A_t=AP$ for a permutation matrix
$P$, then $A(M_t-P)=0$.  Since $A$ has column rank $K$, it has a left inverse,
and therefore $M_t=P$.  The identity is the only permutation matrix in a
sufficiently small neighborhood of $I$.  For nonzero $t$ and nonzero $H$,
$M_t\ne I$, yielding a contradiction.  Finally, the linear constraint
$H\ones=0$ supplies $K$ independent equations on $K^2$ entries, so its
solution space has dimension $K(K-1)$.  Full column rank makes these local
directions distinct at $A$.
\end{proof}

Any strictly positive router can be written as a softmax: for fixed
$\tau>0$, take $z_{ij}=\tau\log A_{ij}$, up to an arbitrary additive constant
in each row.  Thus softmax normalization does not eliminate the orbit.

\paragraph{Entropy-changing gauges.}
For $M_s$ in \eqref{eq:uniformizing_gauge}, the eigenvalue along
$\ones$ is one and every eigenvalue orthogonal to $\ones$ is $1-s$; hence it
is invertible for $0<s<1$.  Since
\begin{equation}
  A\frac{\ones\ones^\top}{K}
  =\frac{\ones_L\ones^\top}{K}
  =U,
\end{equation}
we have $AM_s=(1-s)A+sU$.  Strict concavity of Shannon entropy gives
\begin{equation}
 H((1-s)a_i+su)
 \ge(1-s)H(a_i)+sH(u)>H(a_i)
\end{equation}
for every nonuniform $a_i$, because $u$ is the unique entropy maximizer.
The effective parameters remain exactly unchanged after replacing $B$ by
$M_s^{-1}B$.  As $s\uparrow1$, the router can approach uniform while the
basis norm may grow.  A norm penalty can prefer one member of this orbit, but
that is a gauge-fixing optimization convention rather than identification by
the data.

\paragraph{Verification of the argmax counterexample.}
The matrix $M$ in \eqref{eq:argmax_example} is positive, row
stochastic, and has determinant $0.5$.  Direct multiplication gives
\begin{equation}
 AM=
 \begin{bmatrix}
 0.45&0.55\\
 0.55&0.45\\
 0.60&0.40\\
 0.85&0.15
 \end{bmatrix}.
\end{equation}
Thus two layer assignments flip.  If $B$ has two distinct rows, those rows are
affinely independent; \eqref{eq:app_affine_rank} shows that
$M^{-1}B$ remains so.  The example therefore does not rely on collapsed
bases.

\begin{proof}[Proof of Proposition~\ref{thm:end_to_end_impossibility}]
Let $(A,B)$ and $(A',B')$ be the two factorizations.  Gauge equivalence gives
$AB=A'B'=\Theta$.  By assumption, for every input $x$,
\begin{equation}
  F(x;A,B)=F(x;AB)=F(x;A'B')=F(x;A',B').
\end{equation}
The complete input--output process, including its distribution under any
design over inputs, is therefore identical in the two worlds.  Every
possibly randomized estimator has the same output distribution in both
worlds.  If the two target router summaries differ, the estimator cannot equal
both with probability one.  The argument already grants the estimator all
function values, so increasing the sample size cannot repair the failure.
\end{proof}

The proposition assumes that factors enter the architecture only through the
parameter product.  If all nonlinear basis modules are executed separately
and their outputs are mixed, an arbitrary linear change of basis need not
remain in that model class.  No claim for such output-space mixtures follows
from this proof.

\subsection{Optimizer-response stabilizer}
\label{app:response_proofs}

We use the convention of Section~\ref{sec:response_identifiability}.  Set
\begin{equation}
 \lambda_Z=\eta_Z/\tau^2,
 \qquad
 C_i=\operatorname{Diag}(a_i)-a_ia_i^\top,
 \qquad
 H_i=B^\top C_i^2B.
 \label{eq:app_response_notation}
\end{equation}
With rows of the input and output represented as column vectors, the
\((i,j)\) block of \(\mathcal R_{A,B}\) is
\begin{equation}
 [\mathcal R_{A,B}]_{ij}
 =\eta_B(a_i^\top a_j)I_D
  +\mathbf 1\{i=j\}\lambda_ZH_i.
 \label{eq:app_response_blocks}
\end{equation}

\begin{proof}[Proof of Theorem~\ref{thm:response_stabilizer}]
Because both factorizations have rank \(K\), their column spaces equal the
column space of their common product.  Standard full-rank factorization
algebra therefore gives a unique \(M\in\operatorname{GL}(K)\) satisfying
\begin{equation}
 A'=AM,\qquad B'=M^{-1}B.
 \label{eq:app_response_gauge}
\end{equation}
Both routers have unit row sums, so
\(A(M\ones-\ones)=0\).  Full column rank of \(A\) gives
\begin{equation}
 M\ones=\ones.
 \label{eq:app_response_rowsum}
\end{equation}

Equality of the response on every \(G\) is equality of every block in
\eqref{eq:app_response_blocks}.  For \(i\ne j\), it gives
\begin{equation}
 a_i^\top a_j=(a_i')^\top a_j'.
 \label{eq:app_response_offdiag_gram}
\end{equation}
For a diagonal block, define
\(H_i'=(B')^\top(C_i')^2B'\).  Its equality can be rearranged as
\begin{equation}
 H_i-H_i'=c_iI_D,
 \qquad
 c_i=-\frac{\eta_B}{\lambda_Z}
   \bigl(\|a_i\|_2^2-\|a_i'\|_2^2\bigr)
 \label{eq:app_response_scalar_identity}
\end{equation}
for a scalar \(c_i\).  Every simplex covariance satisfies
\(C_i\succeq0\) and \(C_i\ones=0\); full row rank of \(B\) therefore gives
\(H_i\succeq0\) and \(\operatorname{rank}(H_i)\le K-1<D\).  The same holds for
\(H_i'\).  If \(c_i>0\), then \(H_i=H_i'+c_iI_D\) would be positive
definite; if \(c_i<0\), then \(H_i'=H_i-c_iI_D\) would be positive definite.
Both contradict singularity, so \(c_i=0\).  The remaining scalar part of the
diagonal block then gives
\(\|a_i\|_2^2=\|a_i'\|_2^2\).  Together with
\eqref{eq:app_response_offdiag_gram},
\begin{equation}
 AA^\top=A'A'^\top=AMM^\top A^\top.
\end{equation}
Multiplication by a left inverse of \(A\) and its transpose yields
\begin{equation}
 MM^\top=I_K.
 \label{eq:app_response_orthogonal}
\end{equation}

Since \(H_i=H_i'\), substituting \eqref{eq:app_response_gauge} and cancelling
\(B\) with a right inverse gives
\begin{equation}
 C_i^2=M^{-\top}(C_i')^2M^{-1}.
\end{equation}
By \eqref{eq:app_response_orthogonal}, this is
\begin{equation}
 C_i^2=M(C_i')^2M^\top=(MC_i'M^\top)^2.
\end{equation}
Both \(C_i\) and \(MC_i'M^\top\) are positive semidefinite.  Uniqueness of
the positive-semidefinite square root therefore gives
\begin{equation}
 C_i=MC_i'M^\top.
 \label{eq:app_response_jacobian_conjugacy}
\end{equation}

In column convention, \(a_i'=M^\top a_i\).  Expanding
\eqref{eq:app_response_jacobian_conjugacy}, the two rank-one terms cancel and
leave
\begin{equation}
 \operatorname{Diag}(a_i)
 =M\operatorname{Diag}(a_i')M^\top.
 \label{eq:app_response_diagonal_algebra}
\end{equation}
For distinct \(r,s\), the \((r,s)\) entry says
\begin{equation}
 \sum_{\ell=1}^K M_{r\ell}M_{s\ell}a'_{i\ell}=0
 \qquad\text{for every }i.
\end{equation}
The rows of \(A'\) span \(\RR^K\), hence
\(M_{r\ell}M_{s\ell}=0\) for every \(\ell\).  Thus each column of the
orthogonal \(M\) has exactly one nonzero entry, equal to \(+1\) or \(-1\):
\(M\) is a signed permutation.  \Eqref{eq:app_response_rowsum}
removes the negative signs, so \(M=P\) is a permutation matrix.

Conversely, if \(A'=AP\) and \(B'=P^\top B\), then
\(a_i'=P^\top a_i\), \(C_i'=P^\top C_iP\), the router Gram matrix is
unchanged, and
\begin{equation}
 (B')^\top(C_i')^2B'=B^\top C_i^2B.
\end{equation}
Every block in \eqref{eq:app_response_blocks} is therefore unchanged.  This
proves both directions.
\end{proof}

\subsection{A quantitative inverse for the response}
\label{app:response_quantitative}

The exact stabilizer has a linear inverse on any uniformly interior,
well-conditioned, norm-bounded set.  The following constants are explicit so
that the claim is checkable; they are not intended to be numerically tight.
For a linear map between matrix spaces, \(\|\cdot\|_{F\to F}\) denotes the
operator norm induced by the Frobenius norm.

\begin{theorem}[Quantitative optimizer-response stability]
\label{thm:response_stability}
Let \(K\ge2\), \(L\ge K\), and \(D\ge K\).  Let
\(A,A'\in\RR^{L\times K}\) be row-stochastic and let
\(B,B'\in\RR^{K\times D}\).  Suppose
\begin{align}
 \min_{i,j}\{A_{ij},A'_{ij}\}&\ge\alpha>0,
 \label{eq:app_stability_interior}\\
 \sigma_{\min}(A),\sigma_{\min}(A')&\ge s_A>0,
 &
 \sigma_{\min}(B),\sigma_{\min}(B')&\ge s_B>0,
 \label{eq:app_stability_singular}\\
 \|A\|_F,\|A'\|_F&\le L_A,
 &
 \|B\|_F,\|B'\|_F&\le L_B,
 \label{eq:app_stability_norms}\\
 AB&=A'B'.
 \label{eq:app_stability_product}
\end{align}
Both response operators are defined with the same
\(\tau,\eta_Z,\eta_B>0\).  Put
\(\lambda_Z=\eta_Z/\tau^2\) and
\begin{equation}
 \delta=\|\mathcal R_{A,B}-\mathcal R_{A',B'}\|_{F\to F}.
 \label{eq:app_stability_delta}
\end{equation}
Then
\begin{equation}
 \min_{P\in\mathfrak S_K}
 \left(\|A'-AP\|_F^2+\|B'-P^\top B\|_F^2\right)^{1/2}
 \le C\delta,
 \label{eq:app_stability_conclusion}
\end{equation}
where the following is one finite, fully explicit choice.  First define the
factorization and Gram constants
\begin{equation}
 L_M=\max\{1,L_A/s_A\},\qquad
 c_G=\frac{L}{\eta_Bs_A^2},\qquad
 c_H=\frac{2}{\lambda_Zs_B^2}.
 \label{eq:app_stability_constants_gram}
\end{equation}
Next define the covariance and diagonal-algebra constants
\begin{align}
 \mu_*&=\alpha\left(1+\frac{1}{4L_M^2}\right),&
 c_J&=\frac{L_M^2(c_G+c_H)}{\mu_*},&
 c_D&=\frac{\sqrt{LK}}{s_A}c_J.
 \label{eq:app_stability_constants_cov}
\end{align}
Finally define the row-rounding and global constants
\begin{align}
 c_{\rm row}
 &=\left[Kc_D^2+(c_G+Kc_D^2)^2\right]^{1/2},
 \label{eq:app_stability_constants_row}\\
 \delta_0
 &=\min\left\{1,\frac{1}{2c_G},
                    \frac{1}{2\sqrt K\,c_{\rm row}}\right\},
 \label{eq:app_stability_constants_delta0}\\
 C_{\rm loc}
 &=\sqrt{L_A^2+(L_ML_B)^2}\,\sqrt K\,c_{\rm row},&
 C
 &=\max\left\{C_{\rm loc},
       \frac{2\sqrt{L_A^2+L_B^2}}{\delta_0}\right\}.
 \label{eq:app_stability_constants_global}
\end{align}
\end{theorem}

\begin{proof}
Except where a Frobenius norm is displayed explicitly, all matrix norms in
this proof are spectral norms.  We divide the argument into four steps.

\paragraph{Step 1: recovering the router Gram matrix and an almost-orthogonal
gauge.}
Injecting a vector into input block \(j\) and projecting the output onto block
\(i\) are contractions for the Frobenius-induced operator norm.  Therefore
each block of
\(\mathcal R_{A,B}-\mathcal R_{A',B'}\) has spectral norm at most
\(\delta\).  For \(i\ne j\), \eqref{eq:app_response_blocks} gives
\begin{equation}
 \left|(AA^\top-A'A'^\top)_{ij}\right|
 \le\frac{\delta}{\eta_B}.
 \label{eq:app_stability_offdiag}
\end{equation}
For a diagonal block, write
\begin{equation}
 E_i=t_iI_D+\lambda_Z(H_i-H_i'),\qquad
 t_i=\eta_B(\|a_i\|_2^2-\|a_i'\|_2^2),
 \label{eq:app_stability_diagblock}
\end{equation}
where \(H_i'=(B')^\top(C_i')^2B'\).  Both \(H_i\) and \(H_i'\) are singular
positive-semidefinite matrices: their ranks are at most \(K-1<D\).  If
\(t_i\ge0\), evaluating \(E_i\) on a unit vector in
\(\ker(H_i')\) gives a quadratic form at least \(t_i\).  If \(t_i<0\),
evaluating it on a unit vector in \(\ker(H_i)\) gives a quadratic form at most
\(t_i\).  Since \(\|E_i\|_2\le\delta\), in either case
\begin{equation}
 |t_i|\le\delta,\qquad
 \|H_i-H_i'\|_2\le\frac{2\delta}{\lambda_Z}.
 \label{eq:app_stability_diagseparate}
\end{equation}
The diagonal entries of the router Gram difference consequently satisfy the
same bound as \eqref{eq:app_stability_offdiag}, and hence
\begin{equation}
 \|AA^\top-A'A'^\top\|_F\le\frac{L\delta}{\eta_B}.
 \label{eq:app_stability_gram}
\end{equation}

The common product has rank \(K\), so full-rank factorization algebra gives a
unique \(M\in\operatorname{GL}(K)\) with
\begin{equation}
 A'=AM,\qquad B'=M^{-1}B,\qquad M\ones=\ones.
 \label{eq:app_stability_M}
\end{equation}
Moreover,
\begin{equation}
 M=A^\dagger A',\qquad
 M^{-1}=(A')^\dagger A,\qquad
 \|M\|_2,\|M^{-1}\|_2\le L_A/s_A\le L_M.
 \label{eq:app_stability_M_bounds}
\end{equation}
Multiplying \eqref{eq:app_stability_gram} on the left by \(A^\dagger\)
and on the right by \((A^\dagger)^\top\) yields
\begin{equation}
 \|I_K-MM^\top\|_2\le c_G\delta.
 \label{eq:app_stability_almost_orthogonal}
\end{equation}

\paragraph{Step 2: recovering the simplex covariance matrices.}
Let
\(Q=B^\top(BB^\top)^{-1}\).  Then \(BQ=I_K\) and
\(\|Q\|_2\le1/s_B\).  Substitute \(B'=M^{-1}B\) in \(H_i'\), and multiply
\eqref{eq:app_stability_diagseparate} by \(Q^\top\) and \(Q\), to obtain
\begin{equation}
 \|C_i^2-M^{-\top}(C_i')^2M^{-1}\|_2\le c_H\delta.
 \label{eq:app_stability_cov_squares_raw}
\end{equation}
Set \(Y_i=M^\top C_iM\).  The identity
\begin{align}
 Y_i^2-(C_i')^2
 ={}&M^\top C_i(MM^\top-I_K)C_iM \notag\\
 &+M^\top\{C_i^2-M^{-\top}(C_i')^2M^{-1}\}M
 \label{eq:app_stability_cov_squares}
\end{align}
together with \(\|C_i\|_2\le1\),
\eqref{eq:app_stability_M_bounds},
\eqref{eq:app_stability_almost_orthogonal}, and
\eqref{eq:app_stability_cov_squares_raw} gives
\begin{equation}
 \|Y_i^2-(C_i')^2\|_2
 \le L_M^2(c_G+c_H)\delta.
 \label{eq:app_stability_square_bound}
\end{equation}

Strict interiority now supplies the spectral gap that is unavailable near the
simplex boundary.  If \(a\) is a probability vector with
\(\min_j a_j\ge\alpha\), then
\begin{equation}
 C(a)=\operatorname{Diag}(a)-aa^\top
 \succeq\alpha\Pi_{\ones^\perp}.
 \label{eq:app_stability_cov_gap}
\end{equation}
Indeed, write \(a=\alpha\ones+r\), where \(r\ge0\) and
\(\ones^\top r=1-K\alpha\le1\).  For \(x\perp\ones\), weighted
Cauchy--Schwarz gives
\[
 x^\top C(a)x
 =\alpha\|x\|_2^2+\sum_jr_jx_j^2-\left(\sum_jr_jx_j\right)^2
 \ge\alpha\|x\|_2^2.
\]
Both \(Y_i\) and \(C_i'\) have kernel exactly
\(\operatorname{span}\{\ones\}\): for the former, use
\(M^{-1}\ones=\ones\).  They are symmetric, so they preserve
\(\ones^\perp\).  For a unit \(x\perp\ones\), decompose
\(Mx=c\ones+y\), with \(y\perp\ones\).  Since
\(x=c\ones+M^{-1}y\), orthogonality of \(x\) and
\eqref{eq:app_stability_M_bounds} imply
\begin{equation}
 |c|\sqrt K\le L_M\|y\|_2,\qquad
 1\le |c|\sqrt K+L_M\|y\|_2,
\end{equation}
and therefore \(\|y\|_2\ge1/(2L_M)\).  It follows from
\eqref{eq:app_stability_cov_gap} that, on \(\ones^\perp\),
\begin{equation}
 Y_i\succeq\frac{\alpha}{4L_M^2}I,
 \qquad C_i'\succeq\alpha I.
 \label{eq:app_stability_two_gaps}
\end{equation}

On this subspace put \(D_i=Y_i-C_i'\).  The Sylvester equation
\begin{equation}
 Y_iD_i+D_iC_i'=Y_i^2-(C_i')^2
\end{equation}
has the integral solution
\begin{equation}
 D_i=\int_0^\infty
 e^{-tY_i}\{Y_i^2-(C_i')^2\}e^{-tC_i'}\,dt.
\end{equation}
The two lower bounds in \eqref{eq:app_stability_two_gaps} show that the norm of
this inverse is at most \(1/\mu_*\).  Both sides vanish on
\(\operatorname{span}\{\ones\}\), so \eqref{eq:app_stability_square_bound}
proves on the whole space that
\begin{equation}
 \|M^\top C_iM-C_i'\|_2\le c_J\delta.
 \label{eq:app_stability_cov_linear}
\end{equation}
The shared exact kernel is what permits a linear estimate here; a generic
perturbation bound for square roots would only be H\"older of order one half.

\paragraph{Step 3: approximately normalizing the diagonal algebra.}
Because \(a_i'=M^\top a_i\), expanding the two covariance matrices in
\eqref{eq:app_stability_cov_linear} cancels their rank-one terms and gives
\begin{equation}
 \|M^\top\operatorname{Diag}(a_i)M-\operatorname{Diag}(a_i')\|_2
 \le c_J\delta.
 \label{eq:app_stability_diagonal_algebra}
\end{equation}
Let \(\operatorname{Off}(X)\) set the diagonal of \(X\) to zero, let
\(E_{jj}=e_je_j^\top\), and define
\(X_j=\operatorname{Off}(M^\top E_{jj}M)\).  Taking the off-diagonal part is
contractive in Frobenius norm, while
\(\|X\|_F\le\sqrt K\|X\|_2\).  Thus
\eqref{eq:app_stability_diagonal_algebra} implies
\begin{equation}
 \left\|\sum_{j=1}^K A_{ij}X_j\right\|_F
 \le\sqrt K\,c_J\delta
 \qquad(1\le i\le L).
 \label{eq:app_stability_linear_system}
\end{equation}
Apply row \(j\) of \(A^\dagger\) to this system in the direct-sum Frobenius
space.  That row has Euclidean norm at most \(1/s_A\); Cauchy--Schwarz over
the \(L\) rows gives
\begin{equation}
 \|\operatorname{Off}(M^\top E_{jj}M)\|_F
 \le c_D\delta
 \qquad(1\le j\le K).
 \label{eq:app_stability_off_bound}
\end{equation}

\paragraph{Step 4: rounding the gauge to a permutation.}
Write \(M^\top E_{jj}M=r_jr_j^\top\), where \(r_j^\top\) is row \(j\) of
\(M\), and choose \(q\) maximizing \(|r_{jq}|\).  Suppose first that
\(\delta\le\delta_0\).  From
\eqref{eq:app_stability_almost_orthogonal},
\(\|r_j\|_2^2\ge1/2\), so \(r_{jq}^2\ge1/(2K)\).  The squared Frobenius norm
in \eqref{eq:app_stability_off_bound} contains
\(2r_{jq}^2\sum_{\ell\ne q}r_{j\ell}^2\).  Consequently
\begin{align}
 \sum_{\ell\ne q}r_{j\ell}^2
 &\le Kc_D^2\delta^2,
 \label{eq:app_stability_row_tail}\\
 |r_{jq}^2-1|
 &\le(c_G+Kc_D^2)\delta.
 \label{eq:app_stability_row_peak}
\end{align}
The second inequality also uses \(\delta\le1\).  Equations
\eqref{eq:app_stability_row_tail}--\eqref{eq:app_stability_row_peak} show that
row \(j\) is within \(c_{\rm row}\delta\) of a signed coordinate vector.

Collect the selected signed coordinate rows into a matrix \(S\).  By
\eqref{eq:app_stability_constants_delta0},
\begin{equation}
 \|M-S\|_2\le\|M-S\|_F
 \le\sqrt K\,c_{\rm row}\delta\le\frac12.
 \label{eq:app_stability_rounding}
\end{equation}
If two rows selected the same coordinate, \(S\) would be singular and
\(\sigma_{\min}(M)\le1/2\).  This contradicts
\eqref{eq:app_stability_almost_orthogonal}, which gives
\(\sigma_{\min}(M)\ge1/\sqrt2\).  The selected coordinates are therefore
distinct.  Each row of \(M\) sums to one by \eqref{eq:app_stability_M}.  Its
distance from any negative coordinate vector is consequently at least
\(2/\sqrt K\), whereas its selected-row distance in
\eqref{eq:app_stability_rounding} is at most \(1/(2\sqrt K)\).  Every sign is
positive, so \(S=P\) is a permutation matrix and
\begin{equation}
 \|M-P\|_F\le\sqrt K\,c_{\rm row}\delta.
 \label{eq:app_stability_permutation}
\end{equation}

Now
\[
 \|A'-AP\|_F\le L_A\|M-P\|_F.
\]
Using
\(M^{-1}-P^\top=M^{-1}(P-M)P^\top\) also gives
\[
 \|B'-P^\top B\|_F
 \le L_ML_B\|M-P\|_F.
\]
These two bounds and \eqref{eq:app_stability_permutation} prove
\eqref{eq:app_stability_conclusion} with \(C_{\rm loc}\) whenever
\(\delta\le\delta_0\).  If \(\delta\ge\delta_0\), choosing the identity
permutation and using the parameter-set diameter gives
\[
 \left(\|A'-A\|_F^2+\|B'-B\|_F^2\right)^{1/2}
 \le2\sqrt{L_A^2+L_B^2}
 \le\frac{2\sqrt{L_A^2+L_B^2}}{\delta_0}\delta.
\]
Taking the larger of the local and diameter constants completes the proof.
\end{proof}

\paragraph{Why the quantitative assumptions appear.}
The exact product constraint \eqref{eq:app_stability_product} cannot be
replaced by response closeness alone: \((A,B)\) and \((A,-B)\) have identical
responses but opposite effective matrices.  The two singular-value margins
keep the factorization gauge and the right inverse of \(B\) bounded.  The
interior margin creates the uniform covariance gap in
\eqref{eq:app_stability_cov_gap}; without it, the inverse can become
arbitrarily ill-conditioned near a simplex face.  Norm bounds make the
nonlocal diameter argument finite.  Finally, using the same temperature and
block rates is part of specifying the optimizer metric rather than a
technical normalization.  These assumptions are sufficient and transparent,
not asserted to be jointly minimal.

Here are three explicit degenerating families showing that a uniform linear
constant cannot ignore the interior or either rank margin.  They also make
clear that these are stability failures, not counterexamples to the exact
theorem.  Set \(K=L=D=3\), \(U=\ones\ones^\top/3\), and
\begin{equation}
 Q=\frac13
 \begin{bmatrix}
  2&-1&2\\[-1mm]
  -1&2&2\\[-1mm]
  2&2&-1
 \end{bmatrix}.
 \label{eq:app_stability_fixed_Q}
\end{equation}
Then \(Q\) is a non-permutation orthogonal matrix and \(Q\ones=\ones\).

\emph{Router-rank margin.}
For \(0<\varepsilon\le1/4\), let
\[
 A_\varepsilon=U+\varepsilon(I-U),\quad B_\varepsilon=I,\qquad
 A_\varepsilon'=A_\varepsilon Q,\quad B_\varepsilon'=Q^\top .
\]
Both routers are strictly positive, the basis singular values equal one, and
all norms and entrywise interior margins remain uniformly bounded.  However,
\(\sigma_{\min}(A_\varepsilon)=
\sigma_{\min}(A_\varepsilon')=\varepsilon\).
The products and router Gram matrices agree exactly.  At
\(\varepsilon=0\), every router row is uniform and
\(C_i=(I-U)/3\), which commutes with \(Q\); hence the local response terms
also agree.  Define the permutation-orbit distance \(d_\varepsilon\) by
\[
 d_\varepsilon^2
 =\min_P\left\{
   \|A_\varepsilon'-A_\varepsilon P\|_F^2+
   \|B_\varepsilon'-P^\top B_\varepsilon\|_F^2\right\}.
\]
Polynomial dependence on \(\varepsilon\) gives
\begin{align}
 \|\mathcal R_{A_\varepsilon,B_\varepsilon}
       -\mathcal R_{A_\varepsilon',B_\varepsilon'}\|_{F\to F}
 &=O(\varepsilon),\\
 d_\varepsilon&\longrightarrow \min_P\|Q-P^\top\|_F>0.
\end{align}
Thus no uniform inverse constant survives as \(s_A\downarrow0\).

\emph{Basis-rank margin.}
Keep \(Q\) from \eqref{eq:app_stability_fixed_Q}, put
\[
 A=0.3\ones\ones^\top+0.1I,\quad A'=AQ,\qquad
 B_\varepsilon=\ones e_1^\top+\varepsilon I,\quad
 B_\varepsilon'=Q^\top B_\varepsilon .
\]
The routers have fixed positive entries and fixed full rank; all norms remain
bounded, while \(\sigma_{\min}(B_\varepsilon)\to0\).
Because \(C_i\ones=C_i'\ones=0\), all cross terms with
\(\ones e_1^\top\) vanish and
\[
 B_\varepsilon^\top C_i^2B_\varepsilon=\varepsilon^2C_i^2,\qquad
 (B_\varepsilon')^\top(C_i')^2B_\varepsilon'
 =\varepsilon^2Q(C_i')^2Q^\top .
\]
The Gram parts again agree exactly, so the response difference is
\(O(\varepsilon^2)\).  The parameter distance tends to the strictly positive
distance between \(A'=AQ\) and the finite permutation orbit of the full-rank
\(A\).  Hence the constant must also deteriorate as \(s_B\downarrow0\).

\emph{Interior margin.}
Let
\[
 S=\frac1{\sqrt3}
 \begin{bmatrix}0&-1&1\\1&0&-1\\-1&1&0\end{bmatrix},\qquad
 Q_t=\exp(tS),
 \]
and, for sufficiently small \(t>0\), set
\[
 A_t=t\ones\ones^\top+(1-3t)I,\quad B_t=I,\qquad
 A_t'=A_tQ_t,\quad B_t'=Q_t^\top .
\]
Here \(S^\top=-S\) and \(S\ones=0\), so \(Q_t\) is orthogonal and fixes
\(\ones\).  Moreover
\(A_t'=I+t(\ones\ones^\top-3I+S)+O(t^2)\); its off-diagonal
first-order coefficients are \(1\pm1/\sqrt3>0\).  Thus both routers are
strictly positive for all sufficiently small \(t\), while their minimum entry
is \(\Theta(t)\).  Their smallest singular values are
\(1-3t\), and both basis singular values equal one, so every rank and norm
margin stays uniform.  The products and Gram responses agree exactly.  Each
row tends to a simplex vertex, so
\(C(a_{t,i})=O(t)\) and \(C(a_{t,i}')=O(t)\); consequently the remaining
squared-covariance response difference is \(O(t^2)\). For the first row,
direct expansion gives
$\operatorname{tr}(C(a_{t,1})^2)=10t^2+O(t^3)$ and
$\operatorname{tr}(C(a'_{t,1})^2)=12t^2+O(t^3)$.
Orthogonal conjugation preserves trace, so this block's difference has
spectral norm at least $(2/3)t^2+O(t^3)$ before the fixed positive
$\lambda_Z$ multiplier. The full response gap is therefore $\Theta(t^2)$.
On the other hand,
the identity is the nearest permutation for small \(t\), and
\(\|Q_t-I\|_F=\sqrt2\,t+O(t^2)\), so the parameter-orbit distance is
\(\Theta(t)\).  The ratio of parameter distance to response error therefore
diverges like \(1/t\).  This supplies the claimed simplex-face
ill-conditioning while \(s_A,s_B,L_A,L_B\) remain controlled.

\subsection{Finite structured response tomography}
\label{app:response_tomography}

We now prove Corollary~\ref{cor:response_tomography}.  Write

\[
 Q=AA^\top,
 \qquad
 T_i=\lambda_Z B^\top C_i^2B,
 \qquad
 S=\operatorname{row}(\Theta).
\]
Under the rank hypotheses, $\operatorname{rank}(\Theta)=K$ and

\begin{equation}
 S=\operatorname{row}(B),\qquad
 \operatorname{range}(T_i)\subseteq S,\qquad
 T_i x=0\ \text{for every }x\in S^\perp.
 \label{eq:app_tomo_support}
\end{equation}
Indeed, a left inverse of $A$ gives
$B=A^\dagger\Theta$, while $\Theta=AB$ gives the reverse row-space
inclusion; the remaining claims follow from the two outer factors $B^\top$
and $B$ in $T_i$.

Choose $u_1,\ldots,u_K$ and $v_1,\ldots,v_L$ as in the corollary and
define the rows of the $K$ probes by
\begin{equation}
 g_j^{(1)}=v_j+u_1,\qquad
 g_j^{(r)}=u_r\quad(2\le r\le K).
 \label{eq:app_tomo_probes}
\end{equation}
Let $y_i^{(r)}=[\mathcal R_{A,B}(G^{(r)})]_i$.  Then
\begin{align}
 Q_{ij}
 &=\eta_B^{-1}\langle y_i^{(1)},v_j\rangle,
 \label{eq:app_tomo_Q}\\
 T_i u_1
 &=y_i^{(1)}-\eta_B\sum_{j=1}^LQ_{ij}(v_j+u_1),
 \label{eq:app_tomo_local_first}\\
 T_i u_r
 &=y_i^{(r)}-\eta_B\left(\sum_{j=1}^LQ_{ij}\right)u_r,
 \qquad 2\le r\le K.
 \label{eq:app_tomo_local_later}
\end{align}

\begin{proof}[Proof of Corollary~\ref{cor:response_tomography}]
Project the first response onto $S^\perp$.  By
\eqref{eq:app_tomo_support},
\[
 \Pi_{S^\perp}y_i^{(1)}=\eta_B\sum_{j=1}^LQ_{ij}v_j.
\]
Taking its inner product with $v_j$ gives
\eqref{eq:app_tomo_Q}.  Subtracting the now-known shared term from the
unprojected response leaves
$T_i(v_i+u_1)=T_i u_1$, proving
\eqref{eq:app_tomo_local_first}.  For probe $r\ge2$, the same subtraction
leaves $T_i u_r$, which proves \eqref{eq:app_tomo_local_later}.

For every $x\in\RR^D$, \Eqref{eq:app_tomo_support} gives
\[
 T_i x=\sum_{r=1}^K\langle u_r,x\rangle T_i u_r.
\]
Thus Equations~\ref{eq:app_tomo_Q}--\ref{eq:app_tomo_local_later} recover $Q$
and every complete local block $T_i$, and \eqref{eq:response_blocks}
then determines the whole response operator.  Two exact charts of the same
rank-$K$ product have the same $S$, so the same probes recover both
operators.  Equality on those probes is therefore equivalent to equality of
the complete responses, and Theorem~\ref{thm:response_stabilizer} gives the
permutation conclusion and its converse.  The $Z$-level row shifts remain
unobserved.
\end{proof}

The complement-code solve in \eqref{eq:app_tomo_Q} and the $S$-coordinate
solves use orthonormal systems and have condition number one.  That statement
concerns recovery of $Q,T_1,\ldots,T_L$ from response observations; it does
not remove the rank and simplex-boundary conditioning in the inverse from
these components back to the factor chart.

\begin{proposition}[Deterministic observation-noise propagation]
\label{prop:response_tomography_noise}
Suppose the observed response to probe $r$ is
$Y^{(r)}+E^{(r)}$, where
$\|E^{(r)}\|_F\le\epsilon_r$, and use
Equations~\ref{eq:app_tomo_Q}--\ref{eq:app_tomo_local_later} for reconstruction.
Then
\begin{equation}
 \|\widehat Q-Q\|_F\le\epsilon_1/\eta_B.
 \label{eq:app_tomo_noise_Q}
\end{equation}
If $T^{(S)}$ stacks the $K\times K$ matrices representing every $T_i$
in the $u$ basis, then
\begin{equation}
 \|\widehat T^{(S)}-T^{(S)}\|_F
 \le\left[(1+\sqrt{L+1})^2\epsilon_1^2+
 \sum_{r=2}^K(\epsilon_r+\sqrt L\epsilon_1)^2\right]^{1/2}.
 \label{eq:app_tomo_noise_local}
\end{equation}
\end{proposition}

\begin{proof}
Right multiplication by the orthonormal-row complement-code matrix is
nonexpansive, which proves \eqref{eq:app_tomo_noise_Q}.  The first probe has
$G^{(1)}(G^{(1)})^\top=I_L+\ones\ones^\top$, hence
$\|G^{(1)}\|_2=\sqrt{L+1}$.  The error in the reconstructed first local
column is therefore at most
\[
 \epsilon_1+\eta_B\|\widehat Q-Q\|_F\|G^{(1)}\|_2
 \le(1+\sqrt{L+1})\epsilon_1.
\]
For every later probe, the erroneous shared subtraction contributes at most
$\eta_B\|(\widehat Q-Q)\ones\|_2\le\sqrt L\epsilon_1$, in addition to
$\epsilon_r$.  Projection into the orthonormal $u$ basis is nonexpansive;
squaring and summing these $K$ column bounds proves
\eqref{eq:app_tomo_noise_local}.
\end{proof}

This gives deterministic error radii, not a noisy ``if and only if.''  Two
noisy reconstructions can be separated when their component gap exceeds the
sum of the applicable radii.  Biased, unknown, or optimizer-state-dependent
observation error requires a different model.

\paragraph{Constructive exact factor recovery.}
The diagonal-algebra reduction also yields an explicit reconstruction once
$Q$ and the local blocks are known. Let $U\in\RR^{K\times D}$ contain the
orthonormal row-space basis, factor $Q=XX^\top$ with $X$ full column rank,
and put $C=X^\dagger\Theta U^\top$. There is an orthogonal $O$ with
$X=AO$, so $C=O^\top BU^\top$ is invertible. Writing
$T_i^{(S)}=UT_iU^\top$ and $x_i=X_{i,:}^\top$, define
\[
 D_i=\left[C^{-\top}(T_i^{(S)}/\lambda_Z)C^{-1}\right]^{1/2}
       +x_ix_i^\top
     =O^\top\operatorname{diag}(a_i)O,
\]
where the square root is the PSD one. A weighted sum of these symmetric
matrices has distinct eigenvalues whenever the weighted sums of the columns
of $A$ are distinct. Full column rank makes those columns distinct, so
continuous Gaussian weights satisfy this condition almost surely. If $V$
is the resulting orthogonal eigenvector matrix, the diagonals of $V^\top D_iV$
recover router rows up to one common permutation; eigenvector signs cancel.
The basis then follows from $\widehat B=\widehat A^\dagger\Theta$.
This is a constructive use of established joint-diagonalization algebra
\citep{afsari2008sensitivity,bhaskara2014uniqueness}, not an additional
uniqueness claim. Its numerical implementation symmetrizes observations and
clips negative eigenvalues before taking square roots. Such clipping, and
the conditioning of a chosen weighted eigensystem, require separate analysis
under noise; the exact argument is not a stability guarantee for that estimator.

\begin{proposition}[Generic probes for the structured response]
\label{prop:response_tomography_gaussian}
Let $p=D-K\ge1$, retain the exact product and rank assumptions above, and
draw $q$ effective-gradient matrices independently with iid standard
Gaussian entries.  If
\begin{equation}
 q\ge\max\left\{K,\left\lceil\frac{L}{p}\right\rceil\right\},
 \label{eq:app_tomo_gaussian_count}
\end{equation}
then their full response vectors determine $Q,T_1,\ldots,T_L$, and hence
the complete structured response, with probability one.
\end{proposition}

\begin{proof}
Express each probe in fixed orthonormal coordinates for $S\oplus S^\perp$.
Concatenate the $S^\perp$ input coordinates over all probes into
$X_\perp\in\RR^{L\times qp}$.  The corresponding projected responses obey
$Y_\perp=\eta_BQX_\perp$.  When $qp\ge L$, the Gaussian matrix
$X_\perp$ has row rank $L$ almost surely, so it determines $Q$.  After
subtracting that shared term, the $q$ input coordinates in $S$ for any
fixed layer form a Gaussian $K\times q$ matrix.  It has row rank $K$
almost surely when $q\ge K$, and therefore determines that layer's local
block.  Taking a finite union over layers preserves probability one.
\end{proof}

At the released dimensions $(L,K,D)=(12,4,787{,}712)$, four complete
Gaussian probe/response vectors would satisfy
\eqref{eq:app_tomo_gaussian_count} almost surely if the rank solves were
actually retained and checked.  The older eight-probe artifacts retain norm
comparisons rather than a tomographic reconstruction, so they remain a
fixed-alternative test. The original four-probe artifact instead executes the
explicit design \eqref{eq:app_tomo_probes}, evaluates the analytic full-
dimensional response, and records the reconstruction residuals.  It does not
run four full-model autograd/JVP queries; the analytic formula is validated
separately on the declared $4\times32$ extraction.
The additive audit in Table~\ref{tab:autograd-tomography} measures all four
full-dimensional responses by reverse automatic differentiation followed by
a forward-mode JVP of the packed factor map. It checks reconstruction on
two unseen directions per chart and records deterministic noise-bound
diagnostics. These queries use linear objectives on the packed parameters;
they do not differentiate the language-model task loss or an AdamW update.

This finite observation result is a structured operator-probing corollary,
not an independent central innovation.  In particular,
\citet[Sec.~3.3, Eqs.~(7)--(9)]{novak2022fast} reconstruct a finite NTK from
identity NTK--vector products; generic and structured matrix probing is
developed by \citet{chiu2012matrix}, \citet{halikias2024structured}, and
\citet{amsel2026query}.  Relative to a generic $LD$-dimensional identity
reconstruction, Equations~\ref{eq:app_tomo_probes}--\ref{eq:app_tomo_local_later}
only exploit the present model's known support and shared-block structure to
compress the queries.

Several boundaries are essential.  The packing condition $D-K\ge L$ is
sufficient for the explicit construction, not a lower bound for all finite
designs.  When $K=1$, the single probe recovers the shared response and the
softmax local blocks vanish.  Exact common product is required because both
charts must share $S$; a small numerical product residual alone gives no
row-space perturbation theorem.  The implementation accordingly activates the
common-design conclusion only for charts produced by an explicit exact
invertible gauge, with the floating-point product residual recorded separately.
Known $\eta_B>0$ is used in \eqref{eq:app_tomo_Q}; known $\eta_Z>0$ and
temperature remain part of the chart theorem.  Finally, these interventions
are executable through linear checkpoint objectives
$\mathcal L_G(\Theta)=\langle G,\Theta\rangle$, but they need not occur in a
natural minibatch trajectory.

\begin{proposition}[Finite Gaussian probes for a fixed response alternative]
\label{prop:response_gaussian_probes}
Let \(m=LD\), represent a fixed nonzero response difference by
\(\mathsf D\in\RR^{m\times m}\), and draw independent
\(x_1,\ldots,x_q\sim\mathcal N(0,I_m)\).  For every \(0<t<1\),
\begin{equation}
 \Pr\left[
   \max_{\ell\le q}\|\mathsf D x_\ell\|_2
   <t\|\mathsf D\|_2
 \right]
 \le[2\Phi(t)-1]^q.
 \label{eq:app_response_small_ball}
\end{equation}
Thus failure probability at most \(\beta\) is obtained when
\begin{equation}
 q\ge
 \frac{\log(1/\beta)}{-\log(2\Phi(t)-1)}.
 \label{eq:app_response_probe_count}
\end{equation}
At \(t=1/2\), the coefficient is approximately \(1.04\).
This is a fixed-alternative, absolute-threshold statement.  It neither
calibrates a relative numerical threshold nor uniformly certifies equality
of arbitrary response operators when \(q<m\).
\end{proposition}

\begin{proof}
Let \(v\) be a unit right singular vector for the largest singular value of
\(\mathsf D\).  An SVD gives
\(\|\mathsf D x\|_2\ge\|\mathsf D\|_2|v^\top x|\), and
\(v^\top x\sim\mathcal N(0,1)\).  Independence across probes proves
\eqref{eq:app_response_small_ball}; solving the right-hand side for \(q\)
gives \eqref{eq:app_response_probe_count}.

For the uniform limitation, after any \(q<m\) realized probes there exists a
nonzero linear map whose kernel contains their span, so all its observed
responses vanish.  Conversely, \(m\) noiseless Gaussian probes form an
invertible \(m\times m\) probe matrix almost surely and therefore determine an
arbitrary response operator.  Finally, a reported ratio such as
\(\|\mathsf D x\|_2/\max\{\|\mathcal R x\|_2,\|\mathcal R'x\|_2\}\) has a
random denominator; \eqref{eq:app_response_small_ball} does not calibrate it
without additional bounds on those response norms.
\end{proof}

\begin{corollary}[Pointwise probe detection from inverse stability]
\label{cor:response_inverse_gaussian}
Fix two exact product-equivalent charts in one nondegenerate set from
Corollary~\ref{cor:response_stability}, fix them before drawing the probes,
and let $C$ be that set's common inverse constant.  Define
\[
 d_{\rm orb}=\min_{P\in\mathfrak S_K}
 \left(\|A'-AP\|_F^2+\|B'-P^\top B\|_F^2\right)^{1/2}
\]
and $\mathsf D=\mathcal R_{A,B}-\mathcal R_{A',B'}$.  If
$d_{\rm orb}\ge\Delta>0$, then independent standard-Gaussian probes satisfy,
for every $0<t<1$,
\begin{equation}
 \Pr\left[
   \max_{\ell\le q}\|\mathsf D x_\ell\|_2
   \ge t\Delta/C
 \right]
 \ge 1-[2\Phi(t)-1]^q.
 \label{eq:app_response_inverse_small_ball}
\end{equation}
\end{corollary}

\begin{proof}
The inverse bound gives $\|\mathsf D\|_2\ge d_{\rm orb}/C\ge\Delta/C$.
The complement of the event in
\eqref{eq:app_response_inverse_small_ball} is therefore contained in the
small-ball event of Proposition~\ref{prop:response_gaussian_probes}, which
proves the claim.
\end{proof}

At $t=1/2$ and $q=8$, the upper bound on the miss probability is about
$4.62\times10^{-4}$.  This number does not calibrate the released relative
response ratios: the threshold $t\Delta/C$ is absolute, $C$ may be loose,
the alternative and common nondegenerate set must be fixed before the draws,
and the statement is not simultaneous over probe-selected charts.  If each
observed difference has deterministic error at most $\epsilon$, the
corresponding observable lower threshold is only
$t\Delta/C-\epsilon$ when positive.

\paragraph{Why the hypotheses and target matter.}
The following examples separate exact theorem hypotheses from broader
interpretations.
\begin{itemize}
 \item Equality of the effective model is indispensable.  With \(A'=A\) and
 \(B'=-B\), the router Gram and every quadratic matrix
 \(B^\top C_i^2B\) are unchanged, but \(A'B'=-AB\).  The response alone does
 not identify the represented effective matrix.

 \item Router rank cannot be dropped.  For \(K=D=3\), let every row of
 \(A\) be \(\ones^\top/3\), let \(B=I_3\), and let \(Q\) be a
 non-permutation orthogonal rotation on \(\ones^\perp\) with
 \(Q\ones=\ones\).  Then \(A'=AQ=A\), \(B'=Q^\top\), and both the product
 and response are unchanged.  The two factors are not permutation
 equivalent.

 \item Basis rank cannot be dropped.  Let
 \(K=D=3\),
 \(A=0.3\ones\ones^\top+0.1I_3\), and
 \(B=\ones e_1^\top\).  For a sufficiently small non-permutation orthogonal
 rotation \(Q\) fixing \(\ones\), \(A'=AQ\) remains positive and
 row-stochastic, while \(B'=Q^\top B=B\).  The products and router Gram
 matrices agree, and \(C_iB=0\) makes every local softmax response term zero.

 \item Joint router training is material in general.  If \(\eta_Z=0\) and
 \(K\ge3\), choose any
 positive full-rank router \(A\) close enough to uniform, a full-row-rank
 \(B\), and a small non-permutation orthogonal \(Q\) fixing \(\ones\).
 Then \((AQ,Q^\top B)\) is valid and product equivalent, while the remaining
 response \(\eta_BAA^\top G\) is unchanged.

 \item Equality on one task gradient is not operator equality.  At a
 stationary point the observed gradient is zero for every parameterization,
 regardless of whether their response operators differ.  There is also a
 nonzero blind direction away from stationarity.  For a non-permutation exact
 full-rank pair with common row space $S$, suppose $D>K$ and $L>2K$.
 Choose nonzero $v\in S^\perp$.  Since
 $\operatorname{rank}(AA^\top-A'(A')^\top)\le2K<L$, choose nonzero
 $w$ in its kernel and set the rows of $G$ to $g_j=w_jv$.  Both local
 block differences kill $v$, while the shared difference kills $w$, so
 $(\mathcal R_{A,B}-\mathcal R_{A',B'})(G)=0$ even though the two response
 operators differ by Theorem~\ref{thm:response_stabilizer}.  A passive
 trajectory probes only its realized directions.

 \item Strict positivity is automatic for finite softmax logits.  The exact
 algebra above in fact only needs each router row to lie in the closed simplex:
 positive semidefiniteness and \(C_i\ones=0\) make \(H_i\) singular.  Thus
 the formula-defined response theorem extends to boundary rows when all stated
 factor ranks remain full.  A uniform interior margin is material only for a
 quantitative inverse bound.
\end{itemize}
These are sufficient conditions, not a claim of a jointly minimal assumption
set.  In particular, the proof uses \(\eta_B>0\) to recover the router Gram
matrix, but we have not established that \(\eta_B>0\) is necessary once the
complete collection of local softmax terms is available.

\paragraph{Partial audit of the \(\eta_B=0\) boundary.}
Among the partial subcases closed in the released assumption ledger, we state
two simple ones here without promoting them to separate contributions.  Write
an equivalent chart as
\(A'=AM\), \(B'=M^{-1}B\), with \(M\ones=\ones\).  For a router row written
as a column \(p\), put \(q=M^\top p\).  Equality of the remaining local blocks
reduces, after cancelling the full-row-rank basis, to
\begin{equation}
 C(q)^2=M^\top C(p)^2M.
 \label{eq:app_etaB_zero_local}
\end{equation}
If \(M\ge0\), define
\(\mathsf K=\operatorname{Diag}(p)M\operatorname{Diag}(q)^{-1}\).
Then \(\ones^\top\mathsf K=\ones^\top\), and
\eqref{eq:app_etaB_zero_local} is equivalent to
\[
 (\mathsf K-p\ones^\top)^\top(\mathsf K-p\ones^\top)
 =(I-q\ones^\top)^\top(I-q\ones^\top).
\]
Thus every two columns of \(\mathsf K\) are at distance \(\sqrt2\).
They also lie in the probability simplex, whose diameter \(\sqrt2\) is
attained only by distinct vertices.  Hence \(\mathsf K=P\); using
\(\mathsf Kq=p\) in its definition gives \(M=P\).  One positive router row
therefore rules out every non-permutation nonnegative gauge.

For \(K=2\), no sign restriction is needed.  Up to a candidate column swap,
write
\(M=\left[\begin{smallmatrix}x&1-x\\y&1-y\end{smallmatrix}\right]\)
with \(d=x-y>0\).  A row \((p,1-p)\) maps to first coordinate
\(q=y+dp\), and \eqref{eq:app_etaB_zero_local} becomes
\(q(1-q)=d p(1-p)\).  Full router rank supplies two distinct
\(p_1<p_2\).  Subtracting the two equations gives
\(q_1+q_2=p_1+p_2=:S\); adding them and using
\(q_2-q_1=d(p_2-p_1)=d\Delta\) gives
\[
 0=\frac{1-d}{2}\{S(2-S)+d\Delta^2\}.
\]
The bracket is positive, so \(d=1\), \(y=0\), and \(M=I\), before undoing
the possible column swap.  For signed gauges with \(K\ge3\), even when both
\(A\) and \(AM\) are positive and full rank, compatibility of the rowwise
regular-simplex congruences remains unresolved.  Numerical searches that
collapse rank, approach a simplex face, or satisfy only one row are not
treated as a proof or counterexample.  We therefore retain \(\eta_B>0\) in
the main theorem while making no necessity claim.

The theorem is specific to the named Euclidean metric and known block rates.
Natural gradient, Adam/AdamW or momentum state, finite optimizer steps,
factor-specific penalties, and router auxiliary losses define different
response objects.  Most importantly, \((\Theta,\mathcal R)\) identifies a
present optimizer chart, not the factors or graph that historically generated
the data.

\subsection{Positive recovery results}
\label{app:positive_proofs}

\begin{lemma}[Unique barycentric coordinates]
\label{lem:app_barycentric}
If $b_1,\ldots,b_K$ are affinely independent, then
$aB=a'B$ for $a,a'\in\simplex_K$ implies $a=a'$.
\end{lemma}

\begin{proof}
Set $q=a-a'$.  The simplex constraint gives $q\ones=0$, and equality of the
effective parameters gives $qB=0$.  A nonzero $q$ satisfying both equations
would be a nontrivial affine dependence among the rows of $B$.
\end{proof}

\begin{proof}[Proof of Proposition~\ref{thm:anchor_recovery}]
Every effective parameter lies in the convex hull of the bases:
\(
\{\theta_i\}\subseteq\operatorname{conv}\{b_j\}
\).
Separability supplies the reverse containment at the level of convex hulls,
because each $b_j$ occurs as an anchor row of $\Theta$.  Therefore
\begin{equation}
 \operatorname{conv}\{\theta_1,\ldots,\theta_L\}
 =\operatorname{conv}\{b_1,\ldots,b_K\}.
 \label{eq:app_hull}
\end{equation}
Affine independence makes the right side a $(K-1)$-simplex whose vertices are
exactly the unordered set $\{b_1,\ldots,b_K\}$.  Applying the same argument to
$(A',B')$ shows that its bases are the same vertex set.  Thus
$B'=P^\top B$ for a permutation matrix $P$.  Relative to these identified
vertices, Lemma~\ref{lem:app_barycentric} gives unique coordinates for every
row, hence $A'=AP$.

For the one-hot consequence, $\theta_i=b_{c_i}$.  Affine independence implies
that bases are pairwise distinct, so
$\theta_i=\theta_{i'}$ if and only if $c_i=c_{i'}$.  Full basis usage provides
all anchors and makes the basis set observable.  Without full usage, an unused
basis is unconstrained; with duplicate bases, their assignment labels are
observationally interchangeable.
\end{proof}

Separability must restrict the candidate model class, not merely the planted
truth.  If alternative bases may lie outside the observed hull and need not
appear as anchors, a larger simplex can contain the same $\theta_i$ and assign
them different soft barycentric coordinates.

\begin{proof}[Proof of Theorem~\ref{thm:noisy_recovery}]
If $c_i=c_{i'}$, the triangle inequality gives
\begin{equation}
 \|\widehat\theta_i-\widehat\theta_{i'}\|_2
 =\|e_i-e_{i'}\|_2\le2\varepsilon.
\end{equation}
If $c_i\ne c_{i'}$, the reverse triangle inequality gives
\begin{equation}
 \|\widehat\theta_i-\widehat\theta_{i'}\|_2
 \ge\|b_{c_i}-b_{c_{i'}}\|_2-\|e_i\|_2-\|e_{i'}\|_2
 \ge\gamma-2\varepsilon>2\varepsilon.
\end{equation}
The threshold graph therefore contains all and only within-class edges, so its
connected components recover the partition.  A cluster mean equals its basis
plus an average of error vectors.  By convexity of the norm, that average has
norm at most $\varepsilon$.
\end{proof}

The $4\varepsilon$ condition is a clean sufficient margin, not a necessary
one.  Stable recovery without some separation condition is impossible: two
bases closer than the observation uncertainty can exchange assignments
without changing the observed data appreciably.

\subsection{Finite-sample recovery by shared-input probes}
\label{app:probe_proof}

For a pair $(i,j)$ write
\begin{equation}
 X_s^{ij}=\|Y_i(H_s)-Y_j(H_s)\|_2^2,
 \qquad
 d_{ij}^2=\mathbb E X_s^{ij},
 \qquad
 \widehat d_{ij}^2=n^{-1}\sum_{s=1}^nX_s^{ij}.
\end{equation}
The observation $Y_i$ may include seeded measurement noise.  The population
distances and their gap must be computed under that same observation model.
Noise that differs systematically between layers can close or reverse the
gap for the clean functional partition, in which case the theorem does not
apply; it does not silently redefine the recovery target.

\begin{proof}[Proof of Theorem~\ref{thm:probe_recovery}]
Let $N_p=\binom{L}{2}$.  Substitute $u=\delta/4$ into
\Eqref{eq:probe_bernstein}.  The stated sample-size condition implies
for each pair
\begin{equation}
 \Pr\left(
 |\widehat d_{ij}^2-d_{ij}^2|>\delta/4
 \right)
 \le \frac{\eta}{N_p}.
\end{equation}
A union bound gives, with probability at least $1-\eta$,
\begin{equation}
 \max_{i<j}|\widehat d_{ij}^2-d_{ij}^2|\le\delta/4.
 \label{eq:app_uniform_probe}
\end{equation}
On this event, every within-class pair obeys
\begin{equation}
 \widehat d_{ij}^2\le w+\delta/4<(w+b)/2,
\end{equation}
whereas every between-class pair obeys
\begin{equation}
 \widehat d_{ij}^2\ge b-\delta/4>(w+b)/2.
\end{equation}
Thresholding therefore recovers all pairwise equivalence decisions and hence
the partition.

For agglomerative clustering, \Eqref{eq:app_uniform_probe} implies
that every distance between subsets of one true class is smaller than every
distance between subsets of different true classes, for single, complete, or
average linkage.  While more than $K_f$ clusters remain, at least one true class
is split, and the next merge is necessarily within that class.  Induction
shows that cutting the dendrogram at $K_f$ clusters returns the true
partition.  This class count is a property of the functional equivalence
relation and need not equal the number of router bases.
\end{proof}

\paragraph{A bounded special case.}
If the observations are independent across probe draws $s$ with the same
pairwise expectations, and $0\le X_s^{ij}\le R^2$ almost surely for every
pair, Hoeffding's inequality
gives
\begin{equation}
 \Pr(|\widehat d_{ij}^2-d_{ij}^2|>u)
 \le2\exp(-2nu^2/R^4).
\end{equation}
Thus it is sufficient that
\begin{equation}
 n\ge\frac{8R^4}{\delta^2}
 \log\frac{2\binom{L}{2}}{\eta}.
 \label{eq:app_bounded_probe}
\end{equation}
For unbounded probes, \Eqref{eq:probe_bernstein} must be justified by
the probe distribution, module class, and noise model rather than assumed
silently.  Clipping or normalizing probe responses supplies a directly
checkable bounded version at the cost of changing the target metric.

The released paired sweep uses iid Gaussian inputs and
$Y_{is}=r_i(H_s)+\sigma\epsilon_{is}$ with iid Gaussian $\epsilon_{is}$ and a
fixed absolute $\sigma$.  For its finite MLP and Transformer blocks, LayerNorm
with a positive numerical stabilizer bounds the normalized inputs, and the
remaining fixed affine, GELU, softmax-attention, and residual-update maps give a
bounded clean update.  A bounded vector plus Gaussian noise has a
sub-exponential squared norm, so Bernstein constants exist for each frozen
configuration.  The sweep does not estimate uniform numerical values of
$(\nu,c)$ or the population gap $\delta$; consequently its reported empirical
gap is descriptive and is not presented as an application of the theorem's
population certificate.

\paragraph{What the certificate means.}
The theorem targets equality under the chosen $P_H$ and observation model.
It cannot certify global functional identity outside the support of $P_H$.
Nor does it identify a historical generative graph when two distinct teacher
modules happen to agree on that support.

\subsection{Optimization and initialization effects}
\label{app:dynamics_proofs}

\begin{proof}[Proof of Proposition~\ref{prop:factor_dynamics}]
By the chain rule for $\Theta=AB$,
\begin{equation}
 \nabla_A\mathcal L=GB^\top,
 \qquad
 \nabla_B\mathcal L=A^\top G.
\end{equation}
Euclidean gradient flow therefore satisfies
$\dot A=-GB^\top$ and $\dot B=-A^\top G$.  Differentiating the product gives
\begin{align}
 \dot\Theta
 &=\dot A B+A\dot B\\
 &=-G(B^\top B)-(AA^\top)G.
\end{align}
The right-hand side is not determined by $AB$ alone.  If two factorizations
induce different linear operators on $G$, choosing $G$ outside the null space
of their difference produces different initial effective velocities.
Non-orthogonal gauge transformations generically alter both Gram matrices.
\end{proof}

The previous calculation optimizes $A$ directly and is useful because its
preconditioning is explicit.  The constrained softmax model has the same
coordinate dependence.  In the next derivation vectors are columns, so
$\theta_i=B^\top a_i$ and $g_i=\nabla_{\theta_i}\mathcal L$.

\begin{proposition}[Exact softmax effective dynamics]
\label{prop:app_softmax_dynamics}
Let
\begin{equation}
 a_i=\operatorname{softmax}(z_i/\tau),
 \qquad
 J_i=\frac1\tau
 \left(\operatorname{Diag}(a_i)-a_i a_i^\top\right).
\end{equation}
Without factor-specific regularizers, Euclidean gradient flow in $Z,B$ obeys
\begin{align}
 \dot z_i&=-\eta_ZJ_iBg_i,\\
 \dot B&=-\eta_B\sum_{r=1}^L a_rg_r^\top,\\
 \dot\theta_i
 &=-\eta_ZB^\top J_i^2B\,g_i
   -\eta_B\sum_{r=1}^L\langle a_i,a_r\rangle g_r.
 \label{eq:app_softmax_flow}
\end{align}
\end{proposition}

\begin{proof}
The Jacobian of $a_i$ with respect to $z_i$ is the symmetric matrix $J_i$.
Since $\nabla_{a_i}\mathcal L=Bg_i$, the chain rule gives
$\nabla_{z_i}\mathcal L=J_iBg_i$.  Similarly,
$\nabla_B\mathcal L=\sum_r a_rg_r^\top$.  Including the block learning
rates in the flow therefore gives
\begin{equation}
 \dot a_i=J_i\dot z_i=-\eta_ZJ_i^2Bg_i.
\end{equation}
Substituting $\dot a_i$ and $\dot B$ into
$\dot\theta_i=\dot B^\top a_i+B^\top\dot a_i$ yields
\Eqref{eq:app_softmax_flow}.
\end{proof}

\begin{proof}[Proof of Proposition~\ref{prop:softmax_saturation}]
Let
$C_i=\operatorname{Diag}(a_i)-a_i a_i^\top$, so $J_i=C_i/\tau$.
The matrix $C_i$ is positive semidefinite, hence
\begin{align}
 \|C_i\|_2
 &\le\operatorname{tr}(C_i)
 =1-\|a_i\|_2^2\\
 &\le1-(1-\varepsilon)^2
 \le2\varepsilon.
\end{align}
Using $\nabla_{z_i}\mathcal L=J_iBg_i$ gives
\begin{equation}
 \|\nabla_{z_i}\mathcal L\|_2
 \le\|J_i\|_2\|Bg_i\|_2
 \le\frac{2\varepsilon}{\tau}\|Bg_i\|_2.
\end{equation}
The exact router velocity additionally obeys
\begin{equation}
 \|\dot a_i\|_2
 \le\eta_Z\left(\frac{2\varepsilon}{\tau}\right)^2\|Bg_i\|_2.
\end{equation}
The bound is instantaneous: it accounts for both temperature and the current
probability $1-\varepsilon$, while varying $\tau$ at fixed logits changes
$\varepsilon$ as well.  It does not by itself bound cumulative router motion
or guarantee preservation of an argmax over a training interval.
\end{proof}

\begin{proposition}[Symmetric initialization trap]
\label{prop:app_symmetry_trap}
If every router row is uniform and every basis is initialized identically,
then deterministic gradient flow keeps the router uniform and the bases
identical for all time for which the solution exists.
\end{proposition}

\begin{proof}
If all basis rows equal $b$, then
$Bg_i=(b^\top g_i)\ones$.  Since $J_i\ones=0$, every logit gradient is zero.
Under a uniform router, every basis receives the same gradient
$K^{-1}\sum_i g_i$.  Equality of the bases is therefore invariant, which in
turn keeps every router gradient zero.
\end{proof}

The proposition is not a claim that practical floating-point training always
remains symmetric: independent random initialization, parameter-specific
stochastic perturbations, and implementation nondeterminism can break the
equality. Shared minibatch sampling alone preserves this symmetry when every
sampled loss depends on the factors only through $AB$: the same gradient
equalities hold separately for each minibatch. It identifies the
source of specialization as an explicit or stochastic symmetry breaker rather
than information contained in the initial effective model.

\subsection{Detailed claim boundary}
\label{app:claim_boundary}

The formal results support the following distinctions.
\begin{itemize}
 \item Router entropy, smoothness, bandedness, cosine similarity, and argmax
 graphs are gauge-dependent coordinates.  They require gauge-counterfactual
 or initialization-intervention evidence before receiving a structural
 interpretation.

 \item The effective-weight equality partition is gauge invariant.  It is
 still only a parameter-space partition: neural permutation and scaling
 symmetries can make unequal parameters functionally equivalent.

 \item Shared-input probes target functional equivalence under a named probe
 distribution.  They do not establish equality outside its support and need
 not recover the historical graph that generated a teacher.

 \item A generative graph becomes recoverable under explicit conditions such
 as one-hot routing, full basis usage, affine independence or separation, and
 an observation regime strong enough to distinguish the bases.  Low task loss
 alone supplies none of these conditions.

 \item Affine independence is insufficient for an unknown fully soft
 factorization.  It ensures unique barycentric coordinates only after vertices
 are known; anchors or another valid simplex-identification condition are
 still required.

 \item Factor-specific normalization or regularization may choose a useful
 gauge.  It does not retroactively make that coordinate the data-generating
 graph.  Exact equality constraints can reduce the gauge group but require a
 fresh identifiability analysis.

 \item The dynamics results demonstrate coordinate-dependent preconditioning,
 saturation, and a symmetry trap.  They do not prove a universal optimization
 preference for smooth, monotone, periodic, sparse, or low-entropy routers.

 \item No strict function-class hierarchy, generalization improvement, or
 causal performance benefit follows from the algebraic factorization alone.
 Such claims require separate approximation, statistical, and compute-matched
 arguments.
\end{itemize}
